\centerline{\bf \Huge ДЗ. Круги Эйлера}
\centerline{\bf \Huge }

\begin{flushright}
        \scriptsize{\textbf{Ответы без объяснения решения оцениваются в 0 баллов!}}
\end{flushright}

Каждая задача стоит \textbf{3 балла}, но при этом 1 и 2 задачи \textbf{обязательные}! 

За неправильное выполнение 1 и 2 задачи можно получить \textbf{отрицательное} количество баллов!

\problem{1!} В третьей группе обучаются 30 учеников, а в четвертой - 32 . Причём трое учеников обучаются и в третьей, и в четвертой группе одновременно. Сколько суммарно учеников в третьей и четвертой группе?

\problem{2!} В первой и второй группе суммарно 40 учеников. Причём в первой группе 23 ученика, а во второй - 27 учеников. Сколько учеников занимаются и в первой, и во второй группе одновременно?

\problem{3} Однажды начальник ЭлМШ решил раздать всем ученикам чокопаи и бэбифоксы. Не бэбифоксов раздал 777, а не чокопаев - 333, и еще у него осталось 123 сладости. Сколько сладостей было у начальника до раздачи?

\problem{4}  Однажды начальник ЭлМШ решил раздать всем ученикам чокопаи или бэбифоксы. Не бэбифоксы получили 777 учеников и больше ничего не получили, а не чокопаи - 333 ученика и больше ничего не 
получили. При этом 123 ученика получили сразу обе сладости. Сколько детей получили сладости?

\problem{5} Однажды начальник ЭлМШ решил раздать всем ученикам чокопаи или бэбифоксы. Не бэбифоксы получили 777 учеников, а не чокопаи - 333 ученика и еще 123 ученика получили сразу обе сладости. Сколько детей получили сладости?

\problem{6} В прошлом году в ЭлМШ обучались 50 учеников в трёх различных группах, но на протяжении года некоторые меняли свои группы. В первой и второй группе обучались 10 учеников, в первой и третьей - 11 учеников, во второй и третьей - 12 учеников, причём только 1 ученик обучался сразу в трёх группах. Сколько учеников не переходили в другие группы на протяжении года?